% !TEX root =./ECE444_Practice_main.tex
%
% L23 practice set (Antenna Pattern Lab) -- reading a hand-rotated pattern
% trace: dBc extraction, why a time axis cannot be calibrated to angle, scan
% loss and sidelobe asymmetry, far-field radius, and what a chamber removes.

\fullwidth{\textbf{Learning Objective 3.6: } I can distinguish between array factor
and true antenna pattern and account for element pattern effects.
\\[4pt]\normalfont\small Show your work. A \textbf{Documentation} line at the end
is required (write \textbf{None} if you did not collaborate).}

\begin{questions}

\question \textbf{Reading dBc off a Tracking-tab trace.} A cadet walks an HB100
around the PHASER on a $1\ \m$ arc with the beam held at broadside and records the
Signal vs Time trace. The three lobe peaks read $-18.4\ \text{dBFS}$ (main lobe),
$-31.2\ \text{dBFS}$ (left first sidelobe), and $-29.9\ \text{dBFS}$ (right first
sidelobe).

\begin{parts}

	\begin{minipage}[t]{0.58\linewidth}
		\part Convert both sidelobe readings to dBc.
		\begin{solution}[1.1in]
		A dBc value is the lobe peak referred to the main-lobe peak, so subtract:
		\eq{ L_{\text{left}} &= -31.2 - (-18.4) = -12.8\ \text{dBc} \\
		     L_{\text{right}} &= -29.9 - (-18.4) = -11.5\ \text{dBc} }
		\end{solution}
	\end{minipage}\hfill%
	\begin{minipage}[t]{0.37\linewidth}
		\raggedleft \vspace{12pt}
		$L_{\text{left}}$ = \ansbox[1.25in]{$-12.8$ dBc}
		\\[12pt]
		$L_{\text{right}}$ = \ansbox[1.25in]{$-11.5$ dBc}
	\end{minipage}

	\part The cadet's partner raises \textbf{Rx Gain} by $6\ \db$ and repeats the
	walk. State what happens to the three dBFS readings and what happens to the
	two dBc values, and explain why the two answers differ.
		\begin{solution}[1.3in]
		All three dBFS readings rise by about $6\ \db$ as well, because the gain change
		multiplies every received sample by the same factor. The dBc values do not
		change, because each one is a difference between two readings taken through
		the same gain setting and the common term cancels in the subtraction. This
		is why a pattern is reported as a ratio to its own peak: the result does not
		depend on the receiver's absolute scale, which is not calibrated here anyway.
		\end{solution}

	\part A uniform 8-element array has its first sidelobes at $-13$ dBc. Comment
	on the pair measured above: is this array uniformly illuminated, and what
	accounts for the $1.3\ \db$ spread between the two sides?
		\begin{solution}[1.4in]
		Both readings straddle $-13$ dBc within about a decibel and a half, which is
		consistent with uniform illumination measured in a room. A tapered
		illumination would put the sidelobes at least $17\ \db$ down, and a badly
		calibrated array would show a much larger imbalance. The $1.3\ \db$ spread
		between left and right comes from the measurement rather than the array:
		classroom multipath adds roughly $\pm 1\ \db$ of ripple, and the walking speed
		and radius were not identical on the two halves of the arc. The array factor
		of a symmetric uniform array is symmetric at broadside, so the asymmetry has
		no home in the physics of the array itself.
		\end{solution}

\end{parts}

\newpage

\question \textbf{Why the time axis is not an angle axis.} The same trace runs from
$t = 0.00\ \text{s}$, when the source was at $-90\mydeg$, to $t = 7.40\ \text{s}$,
when it reached $+90\mydeg$. The main-lobe peak lands at $t = 3.10\ \text{s}$. The
beam was commanded to broadside.

\begin{parts}

	\part Explain why the lobe \emph{amplitudes} on this trace are usable while the
	lobe \emph{angles} are not.
		\begin{solution}[1.5in]
		The amplitude recorded at any instant is the pattern evaluated at whatever
		angle the source occupied at that instant, and that value does not depend on
		how fast the source was moving. Walking slowly through a lobe spends more
		samples on it and draws it wider on the plot, but the height each sample
		reaches is unchanged. The angle, in contrast, is never recorded at all. The
		x-axis stores elapsed time, and recovering angle from time requires knowing
		the walking speed at every instant, which nothing in the measurement captures.
		\end{solution}

	\begin{minipage}[t]{0.58\linewidth}
		\part Assume a constant walking speed and convert the main-lobe time to an
		angle. Compare it against the commanded beam angle and state the error.
		\begin{solution}[1.2in]
		A constant speed maps time linearly onto the $180\mydeg$ swept:
		\eq{ \theta &= -90\mydeg + 180\mydeg \lp \frac{3.10}{7.40} \rp \\
		            &= -90\mydeg + 75.4\mydeg = -14.6\mydeg }
		The beam was commanded to $0\mydeg$, so the linear map is in error by
		$14.6\mydeg$.
		\end{solution}
	\end{minipage}\hfill%
	\begin{minipage}[t]{0.37\linewidth}
		\raggedleft \vspace{12pt}
		$\theta$ = \ansbox[1.25in]{$-14.6\mydeg$}
		\\[12pt]
		error = \ansbox[1.25in]{$14.6\mydeg$}
	\end{minipage}

	\begin{minipage}[t]{0.58\linewidth}
		\part The true half-power beamwidth is $13.1\mydeg$. If the hand moves as
		much as $30\%$ faster or slower than its average while crossing the main
		lobe, give the range of beamwidths a constant-speed conversion would report.
		\begin{solution}[1.2in]
		A hand moving $30\%$ fast spends $1/1.30$ of the expected time in the lobe,
		and a constant-speed map converts that shortened interval into a narrower
		angle. A hand moving $30\%$ slow does the opposite:
		\eq{ \theta_{\text{fast}} &= \frac{13.1\mydeg}{1.30} = 10.1\mydeg \\
		     \theta_{\text{slow}} &= \frac{13.1\mydeg}{0.70} = 18.7\mydeg }
		The reported beamwidth spans nearly a factor of two on a pattern that never
		changed.
		\end{solution}
	\end{minipage}\hfill%
	\begin{minipage}[t]{0.37\linewidth}
		\raggedleft \vspace{12pt}
		range = \ansbox[1.4in]{$10.1\mydeg$ to $18.7\mydeg$}
	\end{minipage}

	\part Name the piece of range hardware that removes this error and state what
	it records that the hand walk does not.
		\begin{solution}[1.2in]
		A turntable with a shaft encoder. The turntable carries the antenna through a
		commanded rotation, and the encoder tags every recorded sample with the
		angular position it was taken at, typically to a small fraction of a degree.
		The measurement then stores angle directly instead of inferring it from time,
		so no assumption about rotation speed is needed and a non-uniform rotation
		costs nothing.
		\end{solution}

\end{parts}

\newpage

\question \textbf{Scan loss and sidelobe asymmetry.} The projected-aperture rule
puts the peak power gain of a steered array at $\cos\theta_0$ relative to broadside.

\begin{parts}

	\begin{minipage}[t]{0.58\linewidth}
		\part Compute the predicted peak drop at steer angles of $30\mydeg$,
		$45\mydeg$, and $60\mydeg$.
		\begin{solution}[1.3in]
		The rule applies to power, so the drop in dB is $10\mylog{\cos\theta_0}$:
		\eq{ 30\mydeg&: \ 10\mylog{0.866} = -0.6\ \db \\
		     45\mydeg&: \ 10\mylog{0.707} = -1.5\ \db \\
		     60\mydeg&: \ 10\mylog{0.500} = -3.0\ \db }
		\end{solution}
	\end{minipage}\hfill%
	\begin{minipage}[t]{0.37\linewidth}
		\raggedleft \vspace{12pt}
		$30\mydeg$ = \ansbox[1.1in]{$-0.6\ \db$}
		\\[10pt]
		$45\mydeg$ = \ansbox[1.1in]{$-1.5\ \db$}
		\\[10pt]
		$60\mydeg$ = \ansbox[1.1in]{$-3.0\ \db$}
	\end{minipage}

	\begin{minipage}[t]{0.58\linewidth}
		\part A broadside walk peaks at $-18.4\ \text{dBFS}$ and a $30\mydeg$ walk
		peaks at $-19.3\ \text{dBFS}$. Give the measured scan loss, and state whether
		one pair of walks resolves it given $\pm 0.5\ \db$ of run-to-run scatter.
		\begin{solution}[1.3in]
		\eq{ \Delta = -19.3 - (-18.4) = -0.9\ \db }
		The measured $0.9\ \db$ sits $0.3\ \db$ from the predicted value of $0.6\ \db$ and well inside
		the $\pm 0.5\ \db$ scatter of a hand walk. One pair of runs is therefore
		consistent with the prediction but does not resolve it: the same pair of
		numbers would also be consistent with a scan loss anywhere between about
		$0.4$ and $1.4\ \db$ of scan loss. Averaging several walks at each steer angle narrows the
		scatter enough to make the comparison worth reporting.
		\end{solution}
	\end{minipage}\hfill%
	\begin{minipage}[t]{0.37\linewidth}
		\raggedleft \vspace{12pt}
		$\Delta$ = \ansbox[1.25in]{$-0.9\ \db$}
	\end{minipage}

	\part The measured $30\mydeg$ trace shows first sidelobes of $-12.2$ dBc on the
	low-angle side and $-15.2$ dBc on the high-angle side. Explain why they are
	unequal, given that the array factor alone would make them equal.
		\begin{solution}[1.4in]
		The array factor of a uniformly excited array is a function of
		$\sin\theta - \sin\theta_0$, so its sidelobes are equal in that variable and
		the two first sidelobes of a steered beam carry the same array-factor height.
		The radiated pattern is the array factor multiplied by the element factor,
		and the element factor falls off monotonically away from broadside. The
		high-angle sidelobe near $+59\mydeg$ therefore sits where the element pattern
		is much weaker than at the low-angle sidelobe near $+8\mydeg$, and the
		element factor pulls it down by the difference. Asymmetry in a steered
		pattern is evidence of the element factor, not of a fault in the array.
		\end{solution}

	\begin{minipage}[t]{0.58\linewidth}
		\part Real patch elements follow a power pattern nearer
		$\cos^{1.4}\theta$ than $\cos\theta$. Recompute the $60\mydeg$ scan loss and
		state how much worse it is than the ideal-element bound.
		\begin{solution}[1.2in]
		\eq{ 10\mylog{\cos^{1.4} 60\mydeg} &= 14\mylog{0.500} \\ &= -4.2\ \db }
		That is $1.2\ \db$ worse than the $-3.0\ \db$ ideal-element bound. The
		$\cos\theta$ rule is a best case set by the projected aperture alone; a real
		element's own pattern adds to it, and a link budget that scans to $60\mydeg$
		should carry the larger number.
		\end{solution}
	\end{minipage}\hfill%
	\begin{minipage}[t]{0.37\linewidth}
		\raggedleft \vspace{12pt}
		loss = \ansbox[1.25in]{$-4.2\ \db$}
	\end{minipage}

\end{parts}

\newpage

\question \textbf{How far out is far enough.} The PHASER's 8 elements are spaced
$14\ \text{mm}$ apart, so the aperture spans $D = 98\ \text{mm}$ between the outer
element centers. The HB100 runs at $10.525\ \ghz$.

\begin{parts}

	\begin{minipage}[t]{0.58\linewidth}
		\part Find the far-field distance for this aperture and state whether the
		$1\ \m$ arc used in lab is far enough.
		\begin{solution}[1.3in]
		\eq{ \lambda &= \frac{3\times 10^8}{10.525\times 10^9} = 0.0285\ \m \\
		     r &\ge \frac{2D^2}{\lambda} = \frac{2(0.098)^2}{0.0285} = 0.67\ \m }
		The $1\ \m$ arc is comfortably outside the $0.67\ \m$ boundary, so the pattern has settled
		into its far-field shape everywhere along the walk.
		\end{solution}
	\end{minipage}\hfill%
	\begin{minipage}[t]{0.37\linewidth}
		\raggedleft \vspace{12pt}
		$r$ = \ansbox[1.25in]{$0.67\ \m$}
	\end{minipage}

	\begin{minipage}[t]{0.58\linewidth}
		\part A Ku-band panel with a $0.200\ \m$ aperture is to be measured at
		$15.0\ \ghz$. Find its far-field distance.
		\begin{solution}[1.2in]
		\eq{ \lambda &= \frac{3\times 10^8}{15.0\times 10^9} = 0.0200\ \m \\
		     r &\ge \frac{2(0.200)^2}{0.0200} = 4.00\ \m }
		Doubling the aperture and halving the wavelength pushes the requirement out
		by a factor of eight.
		\end{solution}
	\end{minipage}\hfill%
	\begin{minipage}[t]{0.37\linewidth}
		\raggedleft \vspace{12pt}
		$r$ = \ansbox[1.25in]{$4.00\ \m$}
	\end{minipage}

	\begin{minipage}[t]{0.58\linewidth}
		\part The Ku-band panel is measured on the same $1.0\ \m$ classroom arc anyway.
		Find the phase error across the aperture edge and compare it against the
		$22.5\mydeg$ tolerance behind the far-field criterion.
		\begin{solution}[1.3in]
		The extra path to the aperture edge relative to its center is
		$\Delta \approx D^2/8r$, and the phase error is $360\mydeg\Delta/\lambda$:
		\eq{ \Delta &= \frac{(0.200)^2}{8(1.0)} = 0.005\ \m \\
		     \Delta\phi &= \frac{360\mydeg (0.005)}{0.0200} = 90\mydeg }
		That is four times the $22.5\mydeg$ the criterion allows.
		\end{solution}
	\end{minipage}\hfill%
	\begin{minipage}[t]{0.37\linewidth}
		\raggedleft \vspace{12pt}
		$\Delta\phi$ = \ansbox[1.25in]{$90\mydeg$}
	\end{minipage}

	\part Describe what a $90\mydeg$ edge phase error does to the measured pattern,
	and say which pattern feature it corrupts worst.
		\begin{solution}[1.4in]
		The edges of the aperture arrive out of phase with the center, so the
		measurement no longer sees a plane wave and the aperture is effectively
		defocused. The main lobe broadens and its peak drops, but the damage is worst
		in the nulls and the sidelobes: the cancellation that produces a deep null
		depends on precise phase alignment across the aperture, so the nulls fill in
		and the sidelobes read several decibels high. A sidelobe level measured this
		way is pessimistic, and a null depth measured this way is meaningless.
		\end{solution}

\end{parts}

\newpage

\question \textbf{Classroom against chamber.} The lab trace carries about
$\pm 1\ \db$ of ripple from room reflections adding to the direct path.

\begin{parts}

	\part List which of the four hand-measurement artifacts from the lab an
	anechoic chamber removes, and which one it does not remove by itself.
		\begin{solution}[1.5in]
		A chamber removes the reflection artifact directly: absorber on every surface
		leaves only the direct path, so the ripple disappears. A chamber's fixtures
		also remove the angle error and the amplitude instability, because the
		antenna sits on an encoded turntable and the source is bolted at a fixed
		range rather than carried. The far-field condition is not removed by the
		absorber at all; it is a geometric requirement on the separation distance,
		and a chamber satisfies it only if the chamber is long enough for the
		aperture and frequency being measured. A large antenna in a short chamber has
		the same problem as a large antenna in a classroom.
		\end{solution}

	\begin{minipage}[t]{0.58\linewidth}
		\part A single reflected path of relative amplitude $\rho$ adds to a direct
		path of amplitude $1$. Find the $\rho$ that produces a $+1\ \db$ peak in the
		received amplitude, and express it in dB.
		\begin{solution}[1.3in]
		The two paths add in phase at the ripple peak:
		\eq{ 20\mylog{1+\rho} &= 1.0\ \db \\
		     1+\rho &= 10^{0.05} = 1.122 \\
		     \rho &= 0.122 \ \rightarrow \ -18.3\ \db }
		A reflection only $18\ \db$ below the direct path is enough to move the
		reading by a decibel, which is why room reflections matter at levels that
		sound negligible.
		\end{solution}
	\end{minipage}\hfill%
	\begin{minipage}[t]{0.37\linewidth}
		\raggedleft \vspace{12pt}
		$\rho$ = \ansbox[1.25in]{$0.122 = -18.3\ \db$}
	\end{minipage}

	\begin{minipage}[t]{0.58\linewidth}
		\part The main-lobe peak and the sidelobe peak each carry their own
		$\pm 1\ \db$ of ripple. Give the worst-case uncertainty on a sidelobe level
		reported in dBc, and give the range a true $-13$ dBc sidelobe could read.
		\begin{solution}[1.2in]
		A dBc value is a difference of two readings, and in the worst case the two
		errors have opposite signs and add:
		\eq{ \pm 1\ \db + \pm 1\ \db = \pm 2\ \db }
		A true $-13$ dBc sidelobe can therefore be reported anywhere from $-11$ to
		$-15$ dBc without anything being wrong with the array.
		\end{solution}
	\end{minipage}\hfill%
	\begin{minipage}[t]{0.37\linewidth}
		\raggedleft \vspace{12pt}
		uncertainty = \ansbox[1.1in]{$\pm 2\ \db$}
		\\[12pt]
		range = \ansbox[1.25in]{$-11$ to $-15$ dBc}
	\end{minipage}

	\part A customer requires written proof that this array's sidelobes are below
	$-25$ dBc. State whether the classroom measurement can support that claim, and
	give the two reasons that decide it.
		\begin{solution}[1.4in]
		It cannot. The first reason is the uncertainty computed above: a $\pm 2\ \db$
		window on every dBc reading cannot certify a level that must sit below a
		specific number, because a measured $-26$ dBc could be a true $-24$ dBc. The
		second reason is dynamic range. The electrical sweep's noise floor sits about
		$23\ \db$ below the peak, so a $-25$ dBc sidelobe is buried in the floor and
		there is nothing to read. Both problems are what a chamber and a calibrated
		receiver are for, and the claim needs a range measurement.
		\end{solution}

\end{parts}

\vspace{1.5em}
\noindent\textbf{Documentation:} \rule[-2pt]{0.6\linewidth}{0.4pt}

\end{questions}
